\section{The case of smooth divisors}\label{smooth}


Our goal now is to describe the filtrations introduced in the previous section when $Z$ is a smooth divisor. We will then use this to define
Hodge ideals for arbitrary $\jepPubBold{Q}$-divisors. The key result in the smooth case is the following:

\begin{lemma}\label{lem_decomp2}
Let 
 $$\psi\colon Y=\operatorname{Spec} \jepPubBold{C}[t]\longrightarrow X=\operatorname{Spec} \jepPubBold{C}[x]$$ be
  the map given by $\psi^*(x)=t^{\ell}$. If $Z$ is the divisor on $Y$ defined by
$t$, then
we have
an isomorphism of filtered $\jepPubScript{D}_X$-modules
$$\psi_+\jepPubScript{O}_{Y}(*Z)\simeq\mathop{\textstyle\bigoplus}\limits_{j=0}^{\ell-1}\jepPubScript{M}_j,$$
where 
$\jepPubScript{M}_j\simeq\jepPubScript{D}_X/\jepPubScript{D}_X(\partial_x x - j/\ell)$ 
and $F_k\jepPubScript{M}_j$ is generated over $\jepPubScript{O}_X$ by the classes of $1,\partial_x,\ldots,\partial_x^k$.
Moreover, if we consider on $Y$ the $\mu_{\ell}$-action such that every $\lambda\in\mu_{\ell}$ maps $t$ to $\lambda t$, then
$\jepPubScript{M}_j$ is
the component of $\psi_+\jepPubScript{O}_{Y}(*Z)$
on which every  $\lambda\in\mu_{\ell}$ acts by multiplication with $\lambda^j$.
\end{lemma}

\begin{proof}
As usual, it is easier to do the computation for the filtered right $\jepPubScript{D}$-module $\omega_Y(*Z)$ corresponding to $\jepPubScript{O}_Y(*Z)$. 
Note that this is filtered quasi-isomorphic to the complex
$$A^{\jepPubBullet}:\,\,0\longrightarrow \jepPubScript{D}_Y\xrightarrow{\;\textstyle w\;} \omega_Y(Z)\otimes_{\jepPubScript{O}_Y}\jepPubScript{D}_Y\longrightarrow 0,$$
placed in degrees $-1$ and $0$, where $w(1)=(dt/t)\otimes t\partial_t$; see e.g. \cite[Prop. 3.1]{MP1}. 
Since $\psi$ is finite, the functor $\psi_*$ is exact on quasi-coherent $\jepPubScript{O}_Y$-modules, hence
$\psi_+\omega_Y(*Z)$ is computed by the $0$-th cohomology of the complex
$$B^{\jepPubBullet}=\psi_*(A^{\jepPubBullet}\otimes_{\jepPubScript{D}_Y}\jepPubScript{D}_{Y\to X}),$$
with the obvious induced filtration. The definition of $w$
immediately implies that $w\otimes 1_{\jepPubScript{D}_{Y\to X}}$ is injective. 
Note that $dt/t=(1/\ell)dx/x$ and $t\partial_t=\ell x\partial_x$.

In order to describe the complex $B^{\jepPubBullet}$, note that any element of $B^{-1}$ can be uniquely written as $\sum_{j=0}^{\ell-1}t^jP_j$,
with $P_j\in \jepPubScript{D}_X$. Similarly, any element in $B^0$ can be uniquely written as $\sum_{j=0}^{\ell-1}t^j(dx/x)Q_j$, with $Q_j\in\jepPubScript{D}_X$.
Moreover, if $\tau$ is the differential in $B^{\jepPubBullet}$, then
$$\tau\bigg(\sum_{j=0}^{\ell-1}t^jP_j\bigg)=\sum_{j=0}^{\ell-1}t^j\frac{dx}{x}(x\partial_x+j/\ell)P_j,$$
where we use the fact that $t\partial_tt^j=t^jt\partial_t+jt^j$. In other words, we have have an eigenspace decomposition 
$$B^{\jepPubBullet}\simeq \mathop{\textstyle\bigoplus}\limits_{j=0}^{\ell-1}B^{\jepPubBullet}_j,$$
where $B^{\jepPubBullet}_j$ is identified with the complex
$$0\longrightarrow\jepPubScript{D}_X\longrightarrow \jepPubScript{D}_X\longrightarrow 0,$$
with the differential mapping $P$ to $(x\partial_x+j/\ell)P$.
It follows that $B^{\jepPubBullet}$ is filtered quasi-isomorphic to
$$\mathop{\textstyle\bigoplus}\limits_{j=0}^{\ell-1}\jepPubScript{D}_X/(x\partial_x+j/\ell)\jepPubScript{D}_X,$$
where the filtration on the $j$-th component is such that
$$F_{k-1}\left(\jepPubScript{D}_X/(x\partial_x+j/\ell)\jepPubScript{D}_X\right)$$
is the $\jepPubScript{O}_X$-submodule generated by the classes of $1,\partial_x,\ldots,\partial_x^k$.
Moreover, every $\lambda\in\mu_{\ell}$ acts on the $j^{\rm th}$ factor in the above decomposition
by multiplication with $\lambda^j$. 

The assertion in the lemma now follows immediately from the explicit description
of the equivalence between the categories of left and right $\jepPubScript{D}$-modules on $X={\mathbf A}^1$. 
Indeed, recall that if $\tau$ is the $\jepPubBold{C}$-linear endomorphism of the Weyl algebra 
$\Gamma({\mathbf A}^1,\jepPubScript{D}_{{\mathbf A}^1})$
such that $\tau(PQ)=\tau(Q)\cdot\tau(P)$ for all $P$ and $Q$, and such that $\tau(t)=t$ and $\tau(\partial_t)=-\partial_t$, then the left
$\jepPubScript{D}$-module $N$ corresponding to a right $\jepPubScript{D}$-module $M$ is isomorphic to $M$ itself, with scalar multiplication given via the map $\tau$.
Moreover, for filtered $\jepPubScript{D}$-modules, via this isomorphism $F_kN$ corresponds to $F_{k-1}M$.
In particular, we see that if $M=\jepPubScript{D}_X/P\cdot\jepPubScript{D}_X$, then $N\simeq \jepPubScript{D}_X/\jepPubScript{D}_X\cdot \tau(P)$, and we obtain the statement.
\end{proof}

In what follows, we denote by $\lceil\alpha\rceil$ the smallest integer that is $\geqslant\alpha$. For a $\jepPubBold{Q}$-divisor $D=\sum_{i=1}^ra_iD_i$,
we put $\lceil D\rceil=\sum_{i=1}^r\lceil a_i\rceil D_i$. 

\begin{corollary}\label{smooth_case}
If $h\in \jepPubScript{O}_X(X)$ is nonzero and such that the support $Z$ of ${\rm div}(h)$ is smooth (possibly disconnected), then
for every $\alpha\in\jepPubBold{Q}_{> 0}$ the filtration on 
$\jepPubScript{M}(h^{- \alpha})$ is given by
$$F_k\jepPubScript{M}(h^{-\alpha})=\jepPubScript{O}_X\big((k+1)Z-\lceil D\rceil\big)h^{-\alpha}\quad \text{if}\quad k\geqslant 0,$$
where $D=\alpha\cdot{\rm div}(h)$, 
and $F_k\jepPubScript{M}(h^{-\alpha})=0$ if $k<0$.
\end{corollary}

\begin{proof}
We first reduce to the case when $Z={\rm div}(h)$. We can check the assertion in the proposition locally, hence we may assume that
$Z={\rm div}(g)$, for some $g\in\jepPubScript{O}_X(X)$, and $h=ug^m$, for some $u\in\jepPubScript{O}_X^*(X)$. Furthermore, by Remark~\ref{rem_behavior_etale},
it is enough to prove the assertion after passing to a surjective \'{e}tale cover, hence we may assume that $u=v^m$ for some $v\in\jepPubScript{O}_X^*(X)$.
After replacing $g$ by $vg$, we may thus assume that $h=g^m$. In this case we have an isomorphism
of filtered $\jepPubScript{D}$-modules $\jepPubScript{M}(h^{-\alpha})\simeq\jepPubScript{M}(g^{-m\alpha})$, hence we may and will assume that
${\rm div}(h)=Z$.


We consider the smallest positive integer $\ell$ such that $m:=\ell\alpha\in\jepPubBold{Z}$.
If $\ell=1$, then the assertion follows from the formula for the filtration on $\jepPubScript{O}_X(*Z)$ when $Z$ is smooth; see \cite[Prop. 8.2]{MP1}. Therefore from now on we assume $\ell>1$.

The morphism $h\colon X\to\jepPubBold{A}^1$ is smooth over some open neighborhood of $0$. Using Remark~\ref{rem_behavior_etale}, we see that in order to prove the corollary,
we may assume that $X=\jepPubBold{A}^1$ and $h=x$, the standard coordinate on $\jepPubBold{A}^1$.
Consider the Cartesian diagram 
$$
{\def\labelstyle{\textstyle}\xymatrix{V\ar[r]^{j_0}\ar[d]_{p}&W\ar[d]^{g}\\U\ar[r]_{j}&X,}}
$$
where 
$$j_0\colon V=\operatorname{Spec} \jepPubBold{C}[x,x^{-1},y]/(y^{\ell}-x^{-m})\longrightarrow W=\operatorname{Spec} \jepPubBold{C}[x,z]/(z^{\ell}-x^m),\quad j_0^*(z)=y^{-1}.$$
Let
 $\varphi\colon\widetilde{W}= \operatorname{Spec} \jepPubBold{C}[t] \to W$ be the normalization, given by 
$$\varphi^*(x)=t^{\ell} \,\,\,\, \text{and} \,\,\,\, \varphi^*(z)=t^m.$$ 
(Here we use that $\ell$ and $m$ are relatively prime.) 
Note that $\varphi$ is an isomorphism over $V$, hence we have an open embedding $\iota\colon V\hookrightarrow \widetilde{W}$, with complement the smooth divisor $T$ defined by $t$
(in fact, if $a$ and $b$ are integers such that $am+b\ell=1$, then $\iota^*(t)=y^{-a}x^b$). 
We thus have 
$$j_+p_+\jepPubScript{O}_V\simeq \psi_+\iota_+\jepPubScript{O}_V\simeq\psi_+\jepPubScript{O}_{\widetilde{W}}(*T),$$
where $\psi=g\circ\varphi$. We apply Lemma~\ref{lem_decomp2} for $\psi$. 
Note that $\varphi$ is a $\mu_{\ell}$-equivariant morphism if we let each $\lambda\in\mu_{\ell}$ act on $t$ by multiplication with $\lambda^a$.
By considering the behavior with respect to the
$\mu_{\ell}$-action, we see that
 in the decomposition given by the lemma, we have
$\jepPubScript{M}_j\simeq\jepPubScript{M}(x^{-\alpha})$ if and only if $ja\equiv -1$ (mod $\ell$), that is,
$j\equiv -m$ (mod $\ell$). 

Suppose first that $\alpha<1$, in which case the condition for $j$ is that $j=\ell-m$. As a reality check, note that we indeed have an isomorphism
$$\jepPubScript{D}_X/\jepPubScript{D}_X(\partial_xx -(\ell-m)/\ell)\simeq \jepPubScript{M}(x^{-\alpha})$$
that maps the class of $1$ to $x^{-\alpha}$. 
The formula for the filtration on $\jepPubScript{M}(h^{\alpha})$ now follows from Lemma~\ref{lem_decomp2}.
When $\alpha>1$, we put $m=\lceil\alpha\rceil-1$, and use the fact from Remark \ref{twist_positive}, namely that we have an isomorphism of filtered modules
$$\jepPubScript{M}(x^{-\alpha})\longrightarrow\jepPubScript{M}(x^{-\alpha+m}),\quad gx^{-\alpha}\longmapsto (gx^{-m})x^{-\alpha+m},$$
to reduce the assertion to the case $\alpha\in (0,1)$.
This completes the proof of the corollary. 
\end{proof}




